% !PRINTED none

\begin{center}
\gesperrt{Fünfter Abschnitt.}

\gesperrt{Lineare Transformation.}

\rule{0.3\linewidth}{0.4pt}

\S.\,55.

\gesperrt{Einführung der linearen Transformation.}
\end{center}

Wenn wir in irgend welchen Functionen, die von den $m$ Veränderlichen $x_1, x_2 \ldots x_m$ abhängen, für diese Veränderlichen lineare Ausdrücke einsetzen, die von $m$ neuen Veränderlichen $x_1', x_2' \ldots x_m'$ abhängen
\weq{1}{\begin{aligned}
x_1 &= \alpha_1^{(1)} x_1' + \alpha_1^{(2)} x_2' + \cdots + \alpha_1^{(m)} x_m', \\
x_2 &= \alpha_2^{(1)} x_1' + \alpha_2^{(2)} x_2' + \cdots + \alpha_2^{(m)} x_m', \\
& \cdots\cdots\cdots\cdots\cdots\cdots\cdots\cdots \\
x_m &= \alpha_m^{(1)} x_1' + \alpha_m^{(2)} x_2' + \cdots + \alpha_m^{(m)} x_m',
\end{aligned}}
so nennen wir dies eine \gesperrt{lineare Substitution}, und das Ergebniss eine \gesperrt{lineare Transformation}.

Die Grösse
\weq{2}{r = \begin{vmatrix}
\alpha_1^{(1)}, & \alpha_1^{(2)} & \ldots & \alpha_1^{(m)} \\
\alpha_2^{(1)}, & \alpha_2^{(2)} & \ldots & \alpha_2^{(m)} \\
\multicolumn{4}{c}{\cdot\ \cdot\ \cdot\ \cdot\ \cdot} \\
\alpha_m^{(1)}, & \alpha_m^{(2)} & \ldots & \alpha_m^{(m)}
\end{vmatrix}}
heisst die \gesperrt{Substitutionsdeterminante}. Sie muss immer von Null verschieden sein, da sonst durch (1) eine Abhängigkeit zwischen den Variablen $x_1, x_2 \ldots x_m$ ausgedrückt wäre.

Durch eine lineare Substitution geht jede homogene Function der Variablen $x$ in eine homogene Function desselben Grades der Variablen $x'$ über.
